\documentclass[fleqn]{article}
\usepackage{ctex}  
\usepackage{amsmath}
\usepackage{graphicx} 
\usepackage[T1]{fontenc}
\usepackage{lipsum}
\usepackage{fancyhdr}  % 加载 fancyhdr 包
\usepackage{lastpage}
\usepackage{caption}
\usepackage{float} % 在导言区加载


\begin{document}

\section*{练习一}
\begin{equation}
    equation(A) \;\; \theta'' = -\frac{3g}{2L}cos(\theta) 
\end{equation}
\begin{equation}
    equation(B) \;\;
        \begin{cases}
          x(t)=x_c+t-5cos(\theta) \\
          y(t)=5sin(\theta)
        \end{cases}
\end{equation}
The ground is assumed to be parallel to the tangent of the surface of the Earth.\\
$\theta \equiv$ the angle that the ladder forms with the ground. \\
$\omega \equiv$ the angular velocity of the ladder. \\
$\alpha \equiv$ the angular acceleration of the ladder.

\pagebreak
\numberwithin{equation}{section}
\section*{练习二}
\setcounter{section}{1}
and the angle between the ladder and the floor changes with time in accordance with the equation\\
\begin{equation}
    \theta(t) = arccos (\frac{X(t)}{L})
\end{equation}
where X(t) is given by $X(t) = x_0 + t \cdot V(t)$ is given as $ 1 m/s$. \\
Because $\omega = \lim\limits_{\Delta t \to 0} \frac{\Delta \theta}{\Delta t} = \frac{\mathrm{d}\theta}{\mathrm{d}t}$, the angular velocity of the falling end of the ladder can be determined by \\
\begin{equation}
    \omega(t) = - \frac{V(t)+t\frac{dV}{dt}}{L\sqrt{1-\frac{(x_0+tV(t))^2}{L^2}}}
\end{equation}
Furthermore, since $\alpha \equiv \lim\limits_{\Delta t \to 0} \frac{\Delta \omega}{\Delta t} = \frac{d \omega}{dt}$, the angular acceleration of the falling end of the ladder can be given by \\
\begin{equation}
    \alpha(t) = - \frac{(x_0+tV(t))\left(V(t)+t\frac{dV}{dt}\right)^2 + L^2\left(1-\frac{(x_0+tV(t))^2}{L^2}\right) \left(2\frac{dV}{dt}+t\frac{d^2V}{dt^2}\right)}{L^3\left(1-\frac{\left(x_0+tV(t)\right)^2}{L2}\right)^{3/2}}
\end{equation}
\pagebreak

\section*{练习三}
\renewcommand{\theequation}{\arabic{equation}}
A little more manipulation gives us: \\
\setcounter{equation}{15}
\begin{equation}
    \mathrm{v_x}(t) = \left(\frac{k}{m}t+\frac{1}{\mathrm{v_{0x}}}\right)^{-1}
\end{equation}
This give us the velocity of the ball in the $x$-direction at time $t$. \\
We also want the position at time $t$, which we derive by integrating: \\
\begin{equation}
    p_x(t) = \int v_x(t)dt = \int \left(\frac{k}{m}t+\frac{1}{v_{0x}}\right)^{-1}dt = \frac{m}{k} ln |kv_{0x}t+m|
\end{equation}
In the $y$-direction, equations must take into account both air friction and gravity. We can start the same way as with the $x$-direction:
\begin{equation}
    F_y = g - kv_y^2 = ma_y = m\frac{dv_y}{dt}
\end{equation}
\end{document}

